\noindent\textbf{Elección del procedimiento.} Para una pala-varilla articulada en un extremo, $I_b=\int_0^{L_b}r^2(m_b/L_b)dr=m_bL_b^2/3$. Se suman cuatro palas. El balance de momento da $\tau=I\alpha$, y al integrar desde el reposo $\Omega=\alpha t$.

\begin{enumerate}

\item Usar el eje del extremo, no el del centro de la varilla.

\[\fitmath{\mathrm{Ib}=\frac{\mathrm{mb}\cdot {\mathrm{Lb}}^{2}}{3}}\]

\[\fitmath{\mathrm{Ib}=\frac{(90)\cdot {(7)}^{2}}{3}\approx 1470\,\mathrm{kg\,m^{2}}}\]

\item Las cuatro inercias se suman respecto al mismo eje.

\[\fitmath{I=4\cdot \mathrm{Ib}}\]

\[\fitmath{I=4\cdot (1470)\approx 5880\,\mathrm{kg\,m^{2}}}\]

\item Despejar aceleración angular del par constante.

\[\fitmath{\alpha=\frac{\tau}{I}}\]

\[\fitmath{\alpha=\frac{(12000)}{(5880)}\approx 2.040816\,\mathrm{1/s^{2}}}\]

\item Integrar la aceleración angular constante y despejar el tiempo.

\[\fitmath{t=\frac{\Omega}{\alpha}}\]

\[\fitmath{t=\frac{(30)}{(2.040816327)}\approx 14.7\,\mathrm{s}}\]

\end{enumerate}
